\documentclass[10pt,a4paper]{amsart}
\usepackage[margin=19mm]{geometry}
\usepackage{amsmath,amssymb,mathtools}
\usepackage[scr=rsfs,cal=euler]{mathalfa}
\usepackage{xfrac}
\usepackage[unicode,hidelinks,bookmarks=false]{hyperref}
\newcommand{\ie}{i.e.\ }
\newcommand{\cf}{cf.\ }
\newcommand{\resp}{respectively\ }
\newcommand{\Subsection}{\subsection}
\providecommand{\nobreakdash}{\nobreak\hbox{-}}
\setlength{\emergencystretch}{2em}
\multlinegap0pt
\usepackage{tikz}
\usetikzlibrary{cd}
\usepackage{amssymb}
\usepackage[shortlabels]{enumitem}
\newtheorem{theorem}{Theorem}
\newcommand{\FL}[1]{\operatorname{FL}(#1)}
\title{Finiteness of leaps of modules of integrable derivations of algebras of finite type}
\author{Takuya Miyamoto}
\date{Source: arXiv:2512.00690v4 (CC BY 4.0)}
\begin{document}
\maketitle

\noindent\emph{Source and licence.} Finiteness of leaps of modules of integrable derivations of algebras of finite type. Version arXiv:2512.00690v4 (CC BY 4.0), distributed by arXiv under the Creative Commons Attribution 4.0 International licence, http://creativecommons.org/licenses/by/4.0/, as recorded in the arXiv licence metadata for this exact version and shown on the abstract page https://arxiv.org/abs/2512.00690v4. Full licence notice: \texttt{licenses/arxiv-2512.00690-CC-BY-4.0.txt}.\par\smallskip

\noindent\textit{Editorial scope.} Counted unit: the article's theorem global (source0.tex lines 537--539). The article's complete original proof of it is delivered with it (source0.tex lines 541--574, one \texttt{proof} environment of 34 lines). No mathematics is written, repaired or re-derived: the statement and the proof are the article's own lines, in the article's own notation, and no retained line is substituted or re-flowed.  No further source-local context is needed: every cross-reference of every retained line resolves inside the two delivered spans.  Nothing was omitted from any retained span, so the package carries no omission record.  No retained line contains a citation, so the wrapper carries no bibliography and no bibliography entry.  The retained lines contain the article's own diagram code, which the wrapper typesets with the standard tikz packages; no image file is delivered and the package declares no asset dependency.  Editorial formatting only, in addition: an AMS \texttt{amsart} preamble that restates the article's own declarations and macros used by the retained lines, namely the environments \texttt{theorem} on separate counters, together with the standard packages those lines need.  The retained environments print in this document as Theorem 1 in order of appearance, whereas the article prints the same items with the numbering of its own full document.  The complete native source of the article as distributed in the arXiv e-print for this exact version is bundled unchanged as \texttt{sources/190074/source0.tex}.
\begin{theorem}\label{main thm global}
    Let $X$ be a $k$-scheme essentially of finite type. Then $X$ satisfies $\FL{\tau}$ for a sufficiently large integer $\tau$. In particular, the set of leaps of $X$ is finite.
\end{theorem}

\begin{proof}
    We construct inductively a sequence of pairs $(F_1,i_1),(F_2,i_2),\dots$ such that:
    \begin{enumerate}
        \item $i_1\le i_2\le i_3\le \dots$ are nonnegative integers,
        \item $F_n$ is a sub-$\mathcal{O}_X^{p^{i_n}}$-sheaf of $T_{X/k}^1$ that is coherent via $\mathcal{O}_X\to\mathcal{O}_X^{p^{i_n}}$,
        \item $\operatorname{Ob}_X^{p}\subset\operatorname{Ob}_X^{p^2}\subset \dots \subset F_n$ for every $n$,
        \item $F_1\supset F_2\supset F_3\supset \dots$,
        \item for every $n$, $\operatorname{Supp}(F_n/\operatorname{Ob}_X^{p^{i_n}})\supsetneq \operatorname{Supp}(F_{n+1}/\operatorname{Ob}_X^{p^{i_{n+1}}})$.
    \end{enumerate}
    By (v), this sequence must terminate at some $(F_\tau,i_\tau)$ where we have $$\operatorname{Supp}(F_\tau/\operatorname{Ob}_X^{p^{i_\tau}})=\emptyset,$$ i.e., $F_\tau=\operatorname{Ob}_X^{p^{i_\tau}}$. By (iii), $X$ satisfies $\FL{\tau}$ and in particular the set of leaps of $X$ is finite.

    We set $F_1=T_{X/k}^1,\ i_1=0$. Suppose that we have constructed $(F_n, i_n)$ and we would like to construct $(F_{n+1},i_{n+1})$ satisfying (i)-(v). Let us take a generic point $\xi$ of $\operatorname{Supp}(F_n/\operatorname{Ob}_X^{p^{i_n}})$ (i.e. a point without generalizations). Let $R$ be the local ring of $X$ at $\xi$ and let $Y=\operatorname{Spec}R,\ U:=Y\setminus \{\xi\}$. Then $U$ satisfies $\FL{i_n}$ as $\operatorname{Supp}(F_n|_U/\operatorname{Ob}_U^{p^{i_n}})=\emptyset$. Thus, we can apply the results of the previous section to the local $k$-algebra $R$. We can find an integer $i_{n+1}\ge i_{n}$ such that $Y$ satisfies $\FL{i_{n+1}}$. We define $F_{n+1}$ as the subsheaf of $F_n$ consisting of sections $s\in F_n$ such that $$s|_Y\in \operatorname{Ob}_Y^{p^{i_{n+1}}}.$$
    Conditions (i) and (iv) follow from the construction. We have the following cartesian diagram of $\mathcal{O}_X^{p^{i_{n+1}}}$-sheaves
% https://q.uiver.app/#q=WzAsNCxbMCwxLCJGX24iXSxbMSwxLCJcXGlvdGFfKlxcaW90YV4qRl9uIl0sWzEsMCwiXFxpb3RhXypcXG9wZXJhdG9ybmFtZXtPYn1fWV57aV97bisxfX0iXSxbMCwwLCJGX3tuKzF9Il0sWzAsMV0sWzIsMSwiIiwyLHsic3R5bGUiOnsidGFpbCI6eyJuYW1lIjoiaG9vayIsInNpZGUiOiJ0b3AifX19XSxbMywwLCIiLDAseyJzdHlsZSI6eyJ0YWlsIjp7Im5hbWUiOiJob29rIiwic2lkZSI6InRvcCJ9fX1dLFszLDJdXQ==
\[\begin{tikzcd}
	{F_{n+1}} & {\iota_*\operatorname{Ob}_Y^{i_{n+1}}} \\
	{F_n} & {\iota_*\iota^*F_n}
	\arrow[from=1-1, to=1-2]
	\arrow[hook, from=1-1, to=2-1]
	\arrow[hook, from=1-2, to=2-2]
	\arrow[from=2-1, to=2-2]
\end{tikzcd}\]
    where $\iota: Y\to X$. In other words, we have $F_{n+1}=F_n\times_{\iota_*\iota^*F_n}\iota_*\operatorname{Ob}_Y^{i_{n+1}}$. This implies the quasi-coherence, and hence the coherence of $F_{n+1}$ in the sense stated in (ii). Condition (iii) follows from the fact that $Y$ satisfies $\FL{i_{n+1}}$ and from the construction. The restriction of the above diagram to $Y$ is
    % https://q.uiver.app/#q=WzAsNCxbMCwxLCJcXGlvdGFeKkZfbiJdLFsxLDEsIlxcaW90YV4qXFxpb3RhXypcXGlvdGFeKkZfbiJdLFsxLDAsIlxcaW90YV4qXFxpb3RhXypcXG9wZXJhdG9ybmFtZXtPYn1fWV57aV97bisxfX09XFxvcGVyYXRvcm5hbWV7T2J9X1lee2lfe24rMX19Il0sWzAsMCwiXFxpb3RhXipGX3tuKzF9Il0sWzAsMSwiIiwwLHsibGV2ZWwiOjIsInN0eWxlIjp7ImhlYWQiOnsibmFtZSI6Im5vbmUifX19XSxbMiwxLCIiLDIseyJzdHlsZSI6eyJ0YWlsIjp7Im5hbWUiOiJob29rIiwic2lkZSI6InRvcCJ9fX1dLFszLDAsIiIsMCx7InN0eWxlIjp7InRhaWwiOnsibmFtZSI6Imhvb2siLCJzaWRlIjoidG9wIn19fV0sWzMsMl1d
\[\begin{tikzcd}
	{\iota^*F_{n+1}} & {\iota^*\iota_*\operatorname{Ob}_Y^{i_{n+1}}=\operatorname{Ob}_Y^{i_{n+1}}} \\
	{\iota^*F_n} & {\iota^*\iota_*\iota^*F_n}
	\arrow[from=1-1, to=1-2]
	\arrow[hook, from=1-1, to=2-1]
	\arrow[hook, from=1-2, to=2-2]
	\arrow[equals, from=2-1, to=2-2],
\end{tikzcd}\]
    which is also cartesian. This implies $F_{n+1}|_Y=\iota^*F_{n+1}=\operatorname{Ob}_Y^{i_{n+1}}$ and $\xi \notin \operatorname{Supp}(F_{n+1}/\operatorname{Ob}_X^{p^{i_{n+1}}})$. In particular, (v) follows.
\end{proof}

\end{document}
